\begin{abstract}
Energy-efficient large language model (LLM) inference on GPU clusters is a
multi-objective decision problem whose observations have unequal cost and
fidelity.  The best operating point depends jointly on request geometry,
serving runtime, batching, precision, parallelism, topology, and power policy,
yet a trustworthy label may require an entire node or a multi-node experiment.
We present \system, a bounded hybrid agent that plans evidence collection
rather than estimating joules through free-form language-model reasoning.
Given an inference intent, service-level objectives (SLOs), a target cluster,
and a GPU-hour budget, the agent constructs a feasible configuration pool and
maintains a cross-fidelity belief over energy and performance.  At each step it
selects both a configuration and an evidence level, from historical replay and
serving simulation through single-GPU, intra-node, sparse multi-node, and
target-scale execution.  Our Intent-Conditioned Pareto Information Gain
(\ipig) criterion values an action by its expected reduction in uncertainty
about the high-fidelity SLO-feasible Pareto set per GPU-hour.  Typed tools
perform feasibility checking, simulation, telemetry integration, surrogate
updating, and Pareto analysis; a verification gate prevents extrapolated
predictions from becoming deployment recommendations.  We instantiate the
workflow on \bench\ with interchangeable replay and cluster executors and
define an evaluation spanning cross-scale calibration, frontier recovery per
GPU-hour, verified energy savings, and agent reliability.  \system\ thereby
turns sustainable inference tuning into a reproducible agentic HPC workflow in
which every released configuration is supported by measured target-scale
evidence.
\end{abstract}

\begin{IEEEkeywords}
LLM inference, energy efficiency, agentic AI, multi-fidelity optimization,
Pareto search, GPU clusters, power measurement
\end{IEEEkeywords}
